\section{Discussion and Conclusion}
\label{sec:conclusion}
We define and build social world models through explicit representations of agent mental states, actions, and observations (\protocolname).
Our approach captures complex social dynamics systematically by automatically transforming free-form narratives into \protocolname representations, reducing reporting bias and bridging the gap between raw text and actionable social world models.
We achieve SOTA performance on both third-person social reasoning benchmarks and interactive SOTOPIA-hard evaluations by leveraging structured representations and social world models induced from \protocolname.
\update{We view this work as a first step toward general-purpose social world models for LLMs. Despite these gains, important limitations remain: our approach depends on faithful social parsing and state prediction, and errors in representing implicit norms, private beliefs, or long-horizon dynamics can propagate and mislead downstream reasoning or action. Scaling to deeper interactive settings also raises compute and latency costs, and both parsers and SWMs may inherit biases from underlying LLMs and from culturally contingent social conventions. 
Furthermore, we majorly evaluate our approach on text-based social reasoning tasks. We believe that our approach can be extended to more complex social reasoning tasks, such as embodied and vision-based social reasoning tasks.
Looking ahead, future work should explore improving robustness (e.g., uncertainty-aware parsing and consistency checks), enabling efficient multi-step rollouts for interaction, and training specialized models to generate more accurate SWM efficiently. More comprehensive evaluations that incorporate safety, norm adherence, and human preference will also be essential for deploying social world models in the wild.}
We envision \protocolname as a foundation for building general-purpose social world models that support both social reasoning and interaction across diverse domains.
